\section{Introduction}\label{sec:intro}

\paragraph{ICMS is a VAT.}
ICMS (\emph{Imposto sobre Circula\c{c}\~ao de Mercadorias e Servi\c{c}os}) is levied by each Brazilian state on the sale of goods, electricity, interstate and intermunicipal transport and communications, with credit for the tax paid on inputs.
ICMS is recognised internationally as a value-added tax: an OECD working paper describes it as a state-level VAT \citep{arnold2025}, the OECD classification of taxes places all general taxes on value added under heading 5111 whatever the stage at which they are levied \citep{oecd2026guide}, and peer-reviewed studies call it Brazil's state value-added tax \citep{gomes2024,cabello2020}.
Throughout the paper we write \vat; raw database labels keep their Portuguese names (e.g.\ \texttt{ICMS Declarado}) because they are the values of the data.

\paragraph{Two doors into the enrolled debt.}
In a VAT, the amount due is in principle documented by the invoice trail and self-reported by the firm \citep{pomeranz2015}.
Non-compliance can therefore show up in two ways: tax the firm \emph{declared} and did not pay, and tax that only an audit reveals, which the administration \emph{assesses} in an infringement notice \citep{keen2006}.
The compliance literature separates third-party-reported from self-reported bases \citep{kleven2011}, starting from the classical evasion model \citep{allingham1972,slemrod2007}.
In S\~ao Paulo's enrolled \vat{} debt the two doors are visible in the data, and this paper describes them.
It does not test any of these models.

\paragraph{Research question.}
Do declared and assessed \vat{} debt differ in size, concentration, payment and recovery to the point that compliance and collection studies should treat them separately?

\paragraph{Contribution.}
\discutir{the contribution in the authors' words; suggested scope from the notebook: (i) a full-portfolio description, by certificate and by debtor, of the two origins; (ii) the composition of the largest debts by origin; (iii) a replicable protocol that publishes aggregates only}

\paragraph{Outline.}
Section~\ref{sec:data} describes the data and its known gaps, Section~\ref{sec:method} the measures, Section~\ref{sec:results} the results, and Section~\ref{sec:limits} the open questions and limitations.
